\section{Main fixed-separation result}

The contour identity in \cref{obj:thm:exact-contour-identification} turns identification into a finite geometric problem: find a circle on which the observable residual transform is separated from zero, whose interior contains at least one transform zero, and whose associated empirical transforms are stable enough for estimation. This section gives the fixed-\(\delta\) statistical version of that construction. The key inputs are an explicit transform envelope, a numerical zero-localization radius, and an outcome-weighted transform envelope on the contour domain. Together they support a finite contour bank and the total Borel contour statistic used in the main minimax statement.

The first analytic bound converts the exact Luxemburg convention from \cref{obj:def:model-parameters} into a moment-generating-function envelope. This keeps the constants in the zero-localization step tied to the scale convention used in the law classes.

\begin{lemmav}{lem:luxemburg-mgf-envelope}[Luxemburg MGF envelope]
Let \(P\) be a probability law and let \(\psi>0\). Let \(W\) be a real-valued random variable under \(P\). Suppose that:
\begin{itemize}
\item \textbf{(Regularity.)} \(W\) is measurable and \(W\in L^1(P)\).
\item \textbf{(Centering.)} \(E_P W=0\).
\item \textbf{(Luxemburg envelope.)} Under the Luxemburg sub-Gaussian convention of \cref{obj:def:model-parameters}, \(\exp(W^2/\psi^2)\in L^1(P)\) and
\[
E_P\exp(W^2/\psi^2)\le 2.
\]
\end{itemize}
Then, for every \(t\in\mathbb R\),
\[
E_P e^{tW}\le \exp(4\psi^2t^2).
\]
\end{lemmav}

\Cref{obj:lem:luxemburg-mgf-envelope} supplies the real-axis exponential control used to localize zeros of the treatment-innovation transform. The exact constant is part of the statement, so later radius formulas have a numerical envelope constant.

Cumulant separation then forces the transform geometry to contain a zero inside an explicit disk. Write \(A_k\) for the numerical factor in the following statement and \(R_0\) for the zero-localization radius.

\begin{lemmav}{lem:zero-localization}
Fix a primitive parameter tuple \(p\) and a law \(P\) for \(O=(X,T,Y)\). Assume:
\begin{itemize}
\item \(k=r+1\) with \(r\ge2\);
\item \(\eta=T-g_0(X)\) is centered;
\item \(\|\eta\|_{\psi_2}\le\psi_{\eta}\) under the convention in \cref{obj:ass:eta-subgaussian}; and
\item \(|\kappa_k(\eta)|\ge\delta\) as in \cref{obj:ass:cumulant-separation}.
\end{itemize}
Define
\[
A_k=(2^{k+4}k!)^{1/(k-2)},\qquad
R_0=A_k(\psi_{\eta}^2/\delta)^{1/(k-2)},\qquad
M(z)=E_P[e^{z\eta}].
\]
Then there exists \(z_0\in\mathbb C\) such that \(|z_0|\le R_0\) and \(M(z_0)=0\).
\end{lemmav}

\Cref{obj:lem:zero-localization} is the fixed-separation source of the contour bank. The cumulant lower bound \(\delta\) and the treatment-tail scale \(\psi_{\eta}\) determine a disk that contains a transform zero, while the analytic technology is the same general moment-transform geometry used in zero and random-polynomial arguments \citep{MichelenSahasrabudhe2019,EremenkoFryntov2021,DinhGhoshTranTran2021}. The outer search radius \(R_1\) used below enlarges this localization region for contour selection and empirical control.

The remaining population input bounds the outcome-weighted transform on that outer region. It gives a uniform envelope for the numerator side of the contour ratio, with constants depending only on the parameter record and the outer search radius.

\begin{lemmav}{lem:population-numerator-envelope}[Population numerator envelope]
Let \(p\) be a parameter record as in \cref{obj:def:model-parameters}, with sample size \(n\), primitive constants \(C_{\theta},C_g,C_q,\psi_{\eta},\psi_{\xi}\), and outer search radius \(R_1\) as in \cref{obj:def:contour-bank-handle}. Define
\[
C_G(p)=\left[64\{C_q^4+4C_{\theta}^4\psi_{\eta}^4+4\psi_{\xi}^4\}+2\exp\{16R_1C_g+16R_1^2\psi_{\eta}^2\}\right]^{1/2}.
\]
Let \(m\) be a model with observed-data law \(P\) as in \cref{obj:def:plm-model}. Suppose that \(P\) belongs to the non-Gaussian class \(\mathcal P_{\mathrm{NG},n}\) of \cref{obj:def:non-gaussian-class}, where \(n\) is the sample size in \(p\). Then, for every \(z\in\mathbb C\) with \(|z|\le R_1\), the observable outcome-weighted transform \(G_n(z)=E_P[Ye^{zZ_n}]\) from \cref{obj:def:contour-functional} satisfies
\[
|G_n(z)|\le C_G(p).
\]
\end{lemmav}

The envelope in \cref{obj:lem:population-numerator-envelope} is used on the same contour domain as the denominator transform. It combines the bounded regression ranges with the sub-Gaussian treatment and outcome-noise scales, giving the empirical-transform analysis a deterministic population scale in the sense of standard empirical-process risk arguments \citep{BickelKlaassenRitovWellner1993,VanDerVaart1998,VanDerVaartWellner1996}.

We can now state the fixed-separation estimator result. The statistic \(\widehat\theta_{\mathrm{spec},n}\) is the total Borel contour statistic defined in \cref{obj:def:adaptive-contour-estimator}; the contour bank is fixed by the primitive constants before the law and sample vary. The risk statements apply to the deterministic clipped code pair in force at sample size \(n\). This includes code pairs obtained by conditioning on any external training stage, after which the product-law sample used in the theorem is evaluated with those clipped codes held fixed. The represented-execution clause is a conditional reproducibility statement for that same ordinary Borel statistic.

\begin{theoremv}{thm:adaptive-rootn-minimax}[Fixed-code root-\(n\) minimaxity under fixed cumulant separation]
Fix \(r\ge 2\) and positive constants \(\delta,C_{\theta},C_g,C_q,\psi_{\eta},\psi_{\xi}\).  Then there are constants \(c>0\) and \(C>0\), depending only on these fixed primitive constants, such that the following statements hold.

For every measurable covariate space, let \(p_0\) be a parameter record whose primitive entries \(r,\delta,C_{\theta},C_g,C_q,\psi_{\eta},\psi_{\xi}\) equal the fixed constants above, and fix one experiment-wide record consisting of the certified contour-bank input and the represented range input for \(p_0\).  For every parameter record \(p\) with the same primitive entries as \(p_0\), and for every base model satisfying the usual measurability, conditional-mean, and integrability requirements, write \(n\) for the sample size of \(p\).  Let \(\widehat\theta_{\mathrm{spec},n}\) be the total Borel statistic of \cref{obj:def:adaptive-contour-estimator}, constructed from the transported fixed records and the base treatment-code sequence.  Define the fixed-code non-Gaussian class
\[
\mathcal P_{\mathrm{NG},n}(p;\mathrm{base})
=
\{P:\ P\in\mathcal P_{\mathrm{NG},n}\text{ of }\cref{obj:def:non-gaussian-class},\ 
\bar g_n(P)=\bar g_n(\mathrm{base}),\ 
\bar q_n(P)=\bar q_n(\mathrm{base})\}.
\]
Define the fixed-code JMS ACE class
\[
\mathcal P_{\mathrm{ACE},n}^{\mathrm{JMS}}(p;\mathrm{base})
=
\{P:\ P\in\mathcal P_{\mathrm{ACE},n}^{\mathrm{JMS}}\text{ of }\cref{obj:def:jms-ace-class},\ 
\bar g_n(P)=\bar g_n(\mathrm{base}),\ 
\bar q_n(P)=\bar q_n(\mathrm{base})\}.
\]

\begin{itemize}
\item \textbf{(Measurable statistic.)}  The statistic \(\widehat\theta_{\mathrm{spec},n}\) is measurable.
\item \textbf{(Represented execution.)}  For every compiled bounded spectral adapter satisfying the full canonical build-and-compilation specification for \(p\), the transported fixed records, and the base treatment-code sequence, the represented-data execution realizes the same \(\widehat\theta_{\mathrm{spec},n}\) as in \cref{obj:def:adaptive-contour-estimator}.
\item \textbf{(Class relations.)}  Every law in \(\mathcal P_{\mathrm{ACE},n}^{\mathrm{JMS}}\) belongs to \(\mathcal P_{\mathrm{NG},n}\).  If \(s=r\), then \(\mathcal P_{\mathrm{ACE},n}^{\mathrm{JMS}}\) coincides with the \(L^s(P_X)\)-restricted subclass \(\mathcal P_{\mathrm{NG},n}^{\mathrm{ACE}}\) of \cref{obj:def:ace-comparison-subclass}.
\item \textbf{(Non-Gaussian fixed-code risk.)}  If the sample law is the i.i.d. product law of \cref{obj:ass:iid-sampling}, \(\mathcal P_{\mathrm{NG},n}(p;\mathrm{base})\) is nonempty, \(n\ge 2\), and
\[
\varepsilon_{1,n}\le \{4R_1\exp(2C_gR_1)\}^{-1},
\]
where \(R_1\) is the search radius from the contour-bank construction in \cref{obj:def:adaptive-contour-estimator}, then
\[
\frac{c}{n}
\le
\mathfrak R_n^{\mathrm{NG}}
\le
\sup_{P\in\mathcal P_{\mathrm{NG},n}(p;\mathrm{base})}
E_P\!\left[(\widehat\theta_{\mathrm{spec},n}-\theta_0(P))^2\right]
\le
\frac{C}{n},
\]
and
\[
\sup_{P\in\mathcal P_{\mathrm{NG},n}(p;\mathrm{base})}
Q_{P,1-\gamma}\!\left(|\widehat\theta_{\mathrm{spec},n}-\theta_0(P)|\right)
\le
\sqrt{\frac{C}{\gamma n}},
\]
with generalized lower quantiles as in \cref{obj:def:generalized-quantile}.
\item \textbf{(ACE fixed-code risk.)}  If \(\mathcal P_{\mathrm{ACE},n}^{\mathrm{JMS}}(p;\mathrm{base})\) is nonempty, \(n\ge 2\), and
\[
\varepsilon_{1,n}\le \{4R_1\exp(2C_gR_1)\}^{-1},
\]
then
\[
\frac{c}{n}
\le
\inf_{\widetilde\theta}
\sup_{P\in\mathcal P_{\mathrm{ACE},n}^{\mathrm{JMS}}(p;\mathrm{base})}
E_P\!\left[(\widetilde\theta-\theta_0(P))^2\right]
\le
\sup_{P\in\mathcal P_{\mathrm{ACE},n}^{\mathrm{JMS}}(p;\mathrm{base})}
E_P\!\left[(\widehat\theta_{\mathrm{spec},n}-\theta_0(P))^2\right]
\le
\frac{C}{n},
\]
where the infimum is over measurable real-valued estimators of \(O_1,\ldots,O_n\).  Moreover,
\[
\sup_{P\in\mathcal P_{\mathrm{ACE},n}^{\mathrm{JMS}}(p;\mathrm{base})}
Q_{P,1-\gamma}\!\left(|\widehat\theta_{\mathrm{spec},n}-\theta_0(P)|\right)
\le
\sqrt{\frac{C}{\gamma n}}.
\]
\end{itemize}
\end{theoremv}

\Cref{obj:thm:adaptive-rootn-minimax} is the main fixed-separation risk statement. Under fixed primitive constants and the displayed \(L^1(P_X)\) treatment-code radius condition, the contour statistic is measurable and attains the matched \(n^{-1}\) minimax mean-squared-error order over the fixed-code non-Gaussian class. For \(\gamma\in(1/2,1)\), the same statistic controls the generalized lower \((1-\gamma)\)-quantile of absolute error by order \((\gamma n)^{-1/2}\), with the quantile defined in \cref{obj:def:generalized-quantile}; the ACE fixed-code clause states the parallel \(\Theta(n^{-1})\) MSE result on the common JMS ACE class. The executable correspondence is conditional on the full canonical build-and-compilation specification for the parameter record, transported fixed records, and base treatment-code sequence, while the risk inequalities are ordinary sample-law conclusions for the fixed clipped codes.

The intuition follows the population geometry from the previous section. Cumulant separation supplies at least one transform zero in the radius \(R_0\) disk, and the \(L^1(P_X)\) control on \(\bar g_n-g_0\) keeps the code-residual transform within the stability region used by the contour bank. Empirical transform concentration on the fixed bank, combined with the selected-contour perturbation bound, yields the \(C/n\) mean-squared-error upper bound under fixed primitive constants \citep{BickelKlaassenRitovWellner1993,VanDerVaart1998,ChernozhukovEtAl2018DML,NeweyRobins2018}.

The final result packages two sequence-level consequences. Along non-Gaussian experiments with the same fixed primitive constants, the finite bank can be transported across the sequence. On the bounded-outcome Gaussian comparison class, the target itself is identified by the class restrictions as zero, which gives the corresponding minimax risks directly.

\begin{theoremv}{thm:common-experiment-dichotomy}[Common Experiment Dichotomy]
Fix \(r\ge 2\) and positive constants \(\delta,C_{\theta},C_g,C_q,\psi_{\eta},\psi_{\xi}\). For every measurable covariate space, the following hold.

\begin{itemize}
\item \textbf{(Non-Gaussian sequence.)} Let \(p_0\) have these fixed values of \(r,\delta,C_{\theta},C_g,C_q,\psi_{\eta},\psi_{\xi}\), and fix the primitive records \(\mathfrak p_{\star}\) and \(\mathfrak c_{\theta,\star}\) of \cref{obj:def:contour-bank-handle,obj:def:adaptive-contour-estimator}. Let \((p_j)\) be any parameter sequence with the same fixed values of \(r,\delta,C_{\theta},C_g,C_q,\psi_{\eta},\psi_{\xi}\) as \(p_0\), and with common supplied radii \(\varepsilon_{1,n}\) and \(\varepsilon_{2,n}\). Let \((\widehat g_n,\widehat q_n)\) be any measurable deterministic code pair, and write \((\bar g_n,\bar q_n)\) for the clipped pair. Suppose that \(n_j=p_j.n\to\infty\) and, for all sufficiently large \(j\),
\[
\mathcal P_{\mathrm{NG},n_j}(p_j;\bar g_{n_j},\bar q_{n_j})\neq\varnothing,\qquad
n_j\ge 2,\qquad
\varepsilon_{1,n_j}\le
\Bigl(4R_1(p_j)\exp\{2C_gR_1(p_j)\}\Bigr)^{-1}.
\]
Then there are constants \(c,C>0\) such that, with
\(B\) the translated-dyadic contour bank generated from \(p_0\) and \(\mathfrak p_{\star}\), for all sufficiently large \(j\) the transported primitive bank for \(p_j\) generates the same bank \(B\), the statistic \(\widehat\theta_{\mathrm{spec},n_j}\) of \cref{obj:def:adaptive-contour-estimator} is measurable, every compiled bounded spectral adapter satisfying the full canonical build and compilation specification is represented by the corresponding execution on the same transported primitive records, and
\[
\frac{c}{n_j}\le
\mathfrak R_{n_j}^{\mathrm{NG}}(p_j;\widehat g,\widehat q)
\le
\frac{C}{n_j}.
\]
Moreover, the same statistic witnesses the upper bound:
\[
\sup_{P\in \mathcal P_{\mathrm{NG},n_j}(p_j;\bar g_{n_j},\bar q_{n_j})}
\mathbb E_P\!\left[
\bigl(\widehat\theta_{\mathrm{spec},n_j}-\theta_0(P)\bigr)^2
\right]
\le
\frac{C}{n_j}.
\]

\item \textbf{(Gaussian class.)} For every parameter record \(p\), writing \(n\) for its sample size, every \(P\in\mathcal P_{\mathrm G,n}^{\mathrm{JMS}}\) in \cref{obj:def:gaussian-class} satisfies
\[
\theta_0(P)=0.
\]
For every measurable deterministic code pair \((\widehat g_n,\widehat q_n)\), whenever the clipped-code intersection of \(\mathcal P_{\mathrm G,n}^{\mathrm{JMS}}\) is nonempty,
\[
\mathfrak R_{n}^{\mathrm G}(p;\widehat g,\widehat q)=0,
\qquad
\mathfrak M_{n,1-\gamma}^{\mathrm G}(p;\widehat g,\widehat q)=0.
\]
\end{itemize}
\end{theoremv}

\Cref{obj:thm:common-experiment-dichotomy} records the common-experiment form of the fixed-separation result. The non-Gaussian sequence clause states that one translated-dyadic bank serves all sufficiently large experiments sharing the primitive constants, and that the resulting displayed-parameter minimax risk is of order \(n_j^{-1}\). The Gaussian clause gives a separate diagnostic for the bounded-outcome Gaussian comparison class: its member restrictions force \(\theta_0(P)=0\), so both the mean-squared and generalized-quantile minimax risks equal zero on every nonempty fixed-code intersection.

The fixed-\(\delta\) message is therefore geometric and statistical at once. Cumulant separation creates a zero inside a known disk; treatment-code stability preserves the contour geometry for code residuals; empirical-process control of the transform ratio yields parametric MSE over the fixed-code non-Gaussian class; and the same construction aligns with the ACE comparison class under the class relation stated in \cref{obj:thm:adaptive-rootn-minimax}. The finite-record clauses give the reproducibility interface for the same statistic, with the construction details deferred to the certified-construction appendix.
